% GENERATED FILE. Edit canonical lessons, then run npm run generate:books.
% canonical-source-hash: fnv1a64:d2f15ee9dcd73e4a
% canonical-lessons: SA-C526-vyayamabhumih, SA-C526-dhavakah, SA-C526-ksiprah, SA-C526-langhayati, SA-C526-dalam

\chapter{Training Hard}
\label{ch:sa-training-hard}
\hlchaptermodality{526}

\begin{chapteropening}
\textbf{By the end of this chapter:} \emph{I can talk about practice, speed and leaping.}

Complete the last lesson of chapter 526: I can talk about practice, speed and leaping.
\end{chapteropening}

\section[\texorpdfstring{\sk{व्यायामभूमिः} (vyāyāmabhūmiḥ)}{व्यायामभूमिः (vyāyāmabhūmiḥ)}]{\sk{व्यायामभूमिः} (vyāyāmabhūmiḥ) — a gymnasium; an exercise ground}

\label{lesson:SA-C526-vyayamabhumih}



\begin{quote}
Before the new one: say the Sanskrit for hoists, then the Sanskrit for a challenge.
\end{quote}

\subsection*{You'll want to know: \sk{व्यायामभूमिः}}
\textbf{\sk{व्यायामभूमिः}} — \emph{vyāyāmabhūmiḥ} — \textquotedblleft{}a gymnasium; an exercise ground\textquotedblright{}.

\begin{quote}
\sk{अरुणः} \sk{प्रतिदिनं} \sk{व्यायामभूमिं} \sk{गच्छति।}
\end{quote}

\begin{quote}
\emph{aruṇaḥ pratidinaṃ vyāyāmabhūmiṃ gacchati.}
\end{quote}

\begin{quote}
Arun goes to the exercise ground every day.
\end{quote}

A feminine noun.

\begin{cousinweb}
\textbf{\sk{व्यायामभूमिः}} (\emph{vyāyāmabhūmiḥ}) is built from \emph{vyāyāma}, \textquotedblleft{}exercise\textquotedblright{}, and \emph{bhūmi}, \textquotedblleft{}ground\textquotedblright{}; MW gives \emph{vyāyāmabhūmi} as an exercising ground, a gymnasium (MW).
\end{cousinweb}

\begin{grammarlens}[title={What you've built}]
One more word for describing what you lived and felt.
\end{grammarlens}

\subsection*{Your turn}
\begin{itemize}
\raggedright
  \item \emph{Say it:} \emph{vyāyāmabhūmiḥ}
  \item \emph{Say it:} \emph{vyāyāmabhūmiḥ}, once more
  \item \emph{Say it:} \emph{aruṇaḥ pratidinaṃ vyāyāmabhūmiṃ gacchati.} — Arun goes to the exercise ground every day.
  \item \emph{Recall:} say the Sanskrit for hoists, then the Sanskrit for a challenge, then say \emph{vyāyāmabhūmiḥ} again
\end{itemize}

\begin{tcolorbox}[breakable,colback=teal!4,colframe=teal!35!black,title={Before you move on}]
What does \sk{व्यायामभूमिः} mean? (\textquotedblleft{}A gymnasium; an exercise ground\textquotedblright{}.) Say it once more.
\end{tcolorbox}
\section[\texorpdfstring{\sk{धावकः} (dhāvakaḥ)}{धावकः (dhāvakaḥ)}]{\sk{धावकः} (dhāvakaḥ) — a runner}

\label{lesson:SA-C526-dhavakah}



\begin{quote}
Before the new one: say the Sanskrit for a challenge, then the Sanskrit for a gymnasium.
\end{quote}

\subsection*{You'll want to know: \sk{धावकः}}
\textbf{\sk{धावकः}} — \emph{dhāvakaḥ} — \textquotedblleft{}a runner\textquotedblright{}.

\begin{quote}
\sk{धावकः} \sk{क्षेत्रं} \sk{परितः} \sk{धावति।}
\end{quote}

\begin{quote}
\emph{dhāvakaḥ kṣetraṃ paritaḥ dhāvati.}
\end{quote}

\begin{quote}
The runner runs round the field.
\end{quote}

A masculine noun.

\begin{cousinweb}
\textbf{\sk{धावकः}} (\emph{dhāvakaḥ}) comes from the root √dhāv, \textquotedblleft{}to run\textquotedblright{}, which also gives \emph{dhāvati}, with the agent ending \emph{-aka}; MW gives \emph{dhāvaka} as running, a forerunner (MW).
\end{cousinweb}

\begin{grammarlens}[title={What you've built}]
One more word for describing what you lived and felt.
\end{grammarlens}

\subsection*{Your turn}
\begin{itemize}
\raggedright
  \item \emph{Say it:} \emph{dhāvakaḥ}
  \item \emph{Say it:} \emph{dhāvakaḥ}, once more
  \item \emph{Say it:} \emph{dhāvakaḥ kṣetraṃ paritaḥ dhāvati.} — The runner runs round the field.
  \item \emph{Recall:} say the Sanskrit for a challenge, then the Sanskrit for a gymnasium, then say \emph{dhāvakaḥ} again
\end{itemize}

\begin{tcolorbox}[breakable,colback=teal!4,colframe=teal!35!black,title={Before you move on}]
What does \sk{धावकः} mean? (\textquotedblleft{}A runner\textquotedblright{}.) Say it once more.
\end{tcolorbox}
\section[\texorpdfstring{\sk{क्षिप्रः} (kṣipraḥ)}{क्षिप्रः (kṣipraḥ)}]{\sk{क्षिप्रः} (kṣipraḥ) — quick}

\label{lesson:SA-C526-ksiprah}



\begin{quote}
Before the new one: say the Sanskrit for a gymnasium, then the Sanskrit for a runner.
\end{quote}

\subsection*{You'll want to know: \sk{क्षिप्रः}}
\textbf{\sk{क्षिप्रः}} — \emph{kṣipraḥ} — \textquotedblleft{}quick\textquotedblright{}.

\begin{quote}
\sk{सः} \sk{क्षिप्रः} \sk{धावकः} \sk{अस्ति।}
\end{quote}

\begin{quote}
\emph{saḥ kṣipraḥ dhāvakaḥ asti.}
\end{quote}

\begin{quote}
He is a quick runner.
\end{quote}

\begin{cousinweb}
\textbf{\sk{क्षिप्रः}} (\emph{kṣipraḥ}) comes from the root √kṣip, \textquotedblleft{}to throw\textquotedblright{}, as MW lists it: springing back like a bow, and so quick, speedy, swift (MW).
\end{cousinweb}

\begin{grammarlens}[title={What you've built}]
One more word for describing what you lived and felt.
\end{grammarlens}

\subsection*{Your turn}
\begin{itemize}
\raggedright
  \item \emph{Say it:} \emph{kṣipraḥ}
  \item \emph{Say it:} \emph{kṣipraḥ}, once more
  \item \emph{Say it:} \emph{saḥ kṣipraḥ dhāvakaḥ asti.} — He is a quick runner.
  \item \emph{Recall:} say the Sanskrit for a gymnasium, then the Sanskrit for a runner, then say \emph{kṣipraḥ} again
\end{itemize}

\begin{tcolorbox}[breakable,colback=teal!4,colframe=teal!35!black,title={Before you move on}]
What does \sk{क्षिप्रः} mean? (\textquotedblleft{}Quick\textquotedblright{}.) Say it once more.
\end{tcolorbox}
\section[\texorpdfstring{\sk{लंघयति} (laṅghayati)}{लंघयति (laṅghayati)}]{\sk{लंघयति} (laṅghayati) — leaps over}

\label{lesson:SA-C526-langhayati}



\begin{quote}
Before the new one: say the Sanskrit for a runner, then the Sanskrit for quick.
\end{quote}

\subsection*{You'll want to know: \sk{लंघयति}}
\textbf{\sk{लंघयति}} — \emph{laṅghayati} — \textquotedblleft{}leaps over\textquotedblright{}.

\begin{quote}
\sk{धावकः} \sk{रज्जुं} \sk{लंघयति।}
\end{quote}

\begin{quote}
\emph{dhāvakaḥ rajjuṃ laṅghayati.}
\end{quote}

\begin{quote}
The runner leaps over the rope.
\end{quote}

Past (he …): \sk{सः} \sk{रज्जुं} \sk{लंघितवान्}, \emph{saḥ rajjuṃ laṅghitavān}.

\begin{cousinweb}
\textbf{\sk{लंघयति}} (\emph{laṅghayati}) comes from the root √laṅgh, \textquotedblleft{}to leap over, to go beyond\textquotedblright{}; its causative \emph{laṅghayati} has the same sense, and the past participle is \emph{laṅghita} (MW).
\end{cousinweb}

\begin{grammarlens}[title={What you've built}]
One more word for describing what you lived and felt.
\end{grammarlens}

\subsection*{Your turn}
\begin{itemize}
\raggedright
  \item \emph{Say it:} \emph{laṅghayati}
  \item \emph{Say it:} \emph{laṅghayati}, once more
  \item \emph{Say it:} \emph{dhāvakaḥ rajjuṃ laṅghayati.} — The runner leaps over the rope.
  \item \emph{Recall:} say the Sanskrit for a runner, then the Sanskrit for quick, then say \emph{laṅghayati} again
\end{itemize}

\begin{tcolorbox}[breakable,colback=teal!4,colframe=teal!35!black,title={Before you move on}]
What does \sk{लंघयति} mean? (\textquotedblleft{}Leaps over\textquotedblright{}.) Say it once more.
\end{tcolorbox}
\section[\texorpdfstring{\sk{दलम्} (dalam)}{दलम् (dalam)}]{\sk{दलम्} (dalam) — a team; a party}

\label{lesson:SA-C526-dalam}



\begin{quote}
Before the new one: say the Sanskrit for quick, then the Sanskrit for leaps over.
\end{quote}

\subsection*{You'll want to know: \sk{दलम्}}
\textbf{\sk{दलम्}} — \emph{dalam} — \textquotedblleft{}a team; a party\textquotedblright{}.

\begin{quote}
\sk{अस्माकं} \sk{दले} \sk{दश} \sk{क्रीडकाः} \sk{सन्ति।}
\end{quote}

\begin{quote}
\emph{asmākaṃ dale daśa krīḍakāḥ santi.}
\end{quote}

\begin{quote}
Our team has ten players.
\end{quote}

A neuter noun.

\begin{cousinweb}
\textbf{\sk{दलम्}} (\emph{dalam}) comes from the root √dal, \textquotedblleft{}to split\textquotedblright{}: MW gives \emph{dala} as a piece split off, a part; CDIAL also lists a \emph{dala} \textquotedblleft{}a party, a band\textquotedblright{}, perhaps the same word, perhaps from Dravidian (CDIAL 6215); \textquotedblleft{}a team\textquotedblright{} is the modern sense (MW).
\end{cousinweb}

\begin{grammarlens}[title={What you've built}]
One more word for describing what you lived and felt.
\end{grammarlens}

\subsection*{Your turn}
\begin{itemize}
\raggedright
  \item \emph{Say it:} \emph{dalam}
  \item \emph{Say it:} \emph{dalam}, once more
  \item \emph{Say it:} \emph{asmākaṃ dale daśa krīḍakāḥ santi.} — Our team has ten players.
  \item \emph{Recall:} say the Sanskrit for quick, then the Sanskrit for leaps over, then say \emph{dalam} again
\end{itemize}

\begin{tcolorbox}[breakable,colback=teal!4,colframe=teal!35!black,title={Before you move on}]
What does \sk{दलम्} mean? (\textquotedblleft{}A team; a party\textquotedblright{}.) Say it once more.
\end{tcolorbox}
